\section{Phase, registre dodécaphasique et réalisation exceptionnelle de l'état commun}
\label{sec:testigos-descendentes-ch}

Le théorème de la section~\ref{sec:cierre-cardinal-nativo} tranche la question cardinale à l'intérieur de l'état HMT complet. Les constructions qui suivent ne sont pas un appendice ontologique séparable : elles explicitent des coordonnées corrélées du même état enrichi. La double projection conserve des informations complémentaires : valeur archimédienne dans une branche et incidence, frontière et généalogie dans l'autre.

L'holonomie \(\Gamma_9\) compose neuf actualisations \(U_t\) sur l'état enrichi. Le noyau réduit \(\mathcal K_{\mathrm{ph}}\) est une représentation ultérieure sur des émissions déjà construites, qui conserve une phase et une mémoire réduites :
\[
 U_t\longrightarrow\Gamma_9\longrightarrow\mathcal K_{\mathrm{ph}},
 \qquad
 g(k+9)=g(k),\quad
 q_{\mathrm{mem},k+9}=q_{\mathrm{mem},k}+1,\quad
 \widehat x_{k+9}\neq\widehat x_k.
\]
La phase revient, la mémoire avance et l’état ne se réinitialise pas. En revanche, le
sceau dodécaphasique entier est construit au moyen du lecteur signé et de la
transformée de Hadamard :
\[
 \mathcal S_{12}\;(=R_{\mathrm{sgn}}):
 \mathcal X_{\mathrm{TPK}}^{\mathrm{term,enr}}
 \longrightarrow
 \mathcal U_{12}^{\mathrm{sgn}}:=\operatorname{im}\mathcal S_{12},
 \qquad
 U_{\mathrm{term}}\in\mathcal U_{12}^{\mathrm{sgn}},
 \qquad
 K:=\frac14\mathcal H_{12}U\in\mathbb Z^{12}.
\]
La coordonnée \(U_{\mathrm{term}}\) s’obtient à partir de l’état enrichi au moyen
du registre orienté et des lecteurs de canal :
\[
 \mathcal S_{12}(x_{\mathrm{term}})
 =\Pi_H\!\left(
   \mathcal E_{108}^{90,120}
   (\mathcal R_{12}(x_{\mathrm{term}}))\right)
 =U_{\mathrm{term}}.
\]
La section~\ref{sec:k-registro} développe cette évaluation, l'obtention réversible de \(K\) et le retour transversal aux mêmes vingt-cinq canaux. La sérialisation de \(C_{\mathrm{term}}\) utilisée par les contrôles Python est un point d'entrée de vérification, non un générateur mathématique ni une restriction de la preuve imprimée. Le vecteur \(K\) ne produit à lui seul aucune des deux branches. Avec les représentations corrélées \(P,E,\Phi\), il alimente la fermeture entière
\[
 \begin{gathered}
 Z_m:=P_m+E_m-\Phi_m-K_m,\qquad
 R_{N+1}(Z):=\sum_{r\geq0}Z_{N+1+r}1000^{-r-1},\\
 c_{N+1}^{\infty}:=\left\lfloor R_{N+1}(Z)\right\rfloor,\qquad
 s_m=Z_m+c_{m+1}^{\infty},\qquad
 A_m=s_m\bmod1000,\qquad
 c_m=\frac{s_m-A_m}{1000},\\
 \bigl(P_{1:N},E_{1:N},\Phi_{1:N},K_{1:N};
       c_{N+1}=c_{N+1}^{\infty}\bigr)_{N\geq1}
 \xrightarrow{\ \mathsf T_{\alpha,\infty}\ }
 (\alpha_N)_N
 \longmapsto
 \alpha_{\mathrm{HMT}}:=\varprojlim_N\alpha_N.
 \end{gathered}
\]
La condition \(c_{13}=0\) appartient exclusivement à l'exécution isolée de douze blocs et exprime le rationnel \(\alpha_{12}^{(0)}\). Elle n'est pas la frontière de la flèche infinie : dans celle-ci, chaque horizon reçoit le quotient déterminé par sa queue signée ; en particulier, le premier tour de la section complète a \(c_{13}^{\infty}=-1\). La représentation exceptionnelle conserve un codomaine distinct :
\[
 \begin{gathered}
 B_K=\beta_{\mathrm{arq}}(K)=\{i:K_i\ge729\},\qquad
 s^{(0)}=P+E-\Phi-K,\\
 H^-_\alpha=\eta_{\mathrm{arq}}(s^{(0)})
 =\eta_{\mathrm{exc}}(B_K,P_\pi,H_e,H_\varphi,H_{\mathrm{APP}}),\\
 \widetilde{\mathcal I}_*
 =(B_K\subset H^-_\alpha;P_\pi,H_e,H_\varphi,H_{\mathrm{APP}})
 \longmapsto o_*=1
 \longmapsto v_\alpha=4\rho_1+\rho_2+\cdots+\rho_{12},\\
 N_{\alpha,0}
 =\{x\in N(C_W):\langle x,v_\alpha\rangle\equiv0\pmod3\},\qquad
 N_\alpha=N_{\alpha,0}+\mathbb Z\frac{v_\alpha}{3}\cong\Lambda_{24},\\
 C_W\longmapsto
 N(C_W):=\bigcup_{c\in C_W}\bigl(A_2^{12}+\phi(c)\bigr)
 \xrightarrow{\ \operatorname{Nbr}_{v_\alpha}\ }
 N_\alpha\cong\Lambda_{24}
 \longrightarrow V^\natural
 \xrightarrow{\ \operatorname{Aut}\ }\mathbb M.
 \end{gathered}
\]
Le scalaire \(\alpha_{\mathrm{HMT}}\), la face \(B_K\), l’hexade
\(H^-_\alpha\), le vecteur \(v_\alpha\) et le réseau \(N_\alpha\) appartiennent
à des codomaines différents. Ils partagent une provenance ; ce ne sont pas des objets interchangeables.

\begin{proposition}[Unité causale des représentations corrélées]
\label{prop:posterioridad-testigos-ch}
Les objets
\[
 \mathcal K_{\mathrm{ph}},\quad K,\quad
 \pi_{\mathrm{HMT}},\varphi_{\mathrm{HMT}},e_{\mathrm{HMT}},
 \alpha_{\mathrm{HMT}},\quad H^-_\alpha,\quad v_\alpha,\quad
 N_\alpha,\quad V^\natural
\]
ne sont ni des données externes ni des sélecteurs rétroactifs de l'inférence cardinale. Tous sont produits comme coordonnées, représentations ou réalisations des deux projections de l'unique état enrichi. La comparaison cardinale ne consulte pas leurs valeurs conventionnelles comme des oracles, mais elle opère bien sur la structure HMT complète qui conserve leur généalogie, leurs régions et leur mémoire.
\end{proposition}

\begin{proof}
La démonstration cardinale utilise le système inverse d'historiques, la droite intégrale, la clôture sémantique, le classificateur ordinal et les deux injections dans le même état APP--TRIT--TPK. Aucun chiffre conventionnel ne détermine ces applications. Dans l'ordre constructif interne, \(\mathcal K_{\mathrm{ph}}\) factorise la mémoire de \(\Gamma_9\) ; le sceau \(K\) est obtenu réversiblement à partir de la coordonnée terminale produite par \(\mathcal E_{108}^{90,120}\circ\mathcal R_{12}\) ; la quadrirelation s'exprime à partir des sections coinductives ; et le drapeau, le vecteur \(v_\alpha\), le voisin de Leech et \(V^\natural\) sont obtenus par la projection exceptionnelle. Les sections suivantes fournissent les constructions complètes et leurs preuves.
\end{proof}

Sont préservées, en particulier, les séparations de type
\[
 \mathcal K_{\mathrm{ph}}\neq K,\qquad
 \mathcal H_{12}\neq H^-_\alpha,\qquad
 \alpha_{\mathrm{HMT}}\neq v_\alpha\neq N_\alpha.
\]
Il n'existe pas non plus de flèche canonique \(W_{24}\to\Lambda_{24}\) : le plan combinatoire de Witt et le voisin ternaire de Leech sont des prolongements distincts du code commun. Le groupe Monstre agit sur \(V^\natural\), et non sur des chiffres décimaux, des chemins TPK ou des codes de \(\alpha\).
